\lesson{4}{Sep 22 2021 Wen (7:15:54)}{Humor}{Unit 1}

\subsubsection*{vocabulary}

\begin{definition}[Hyperbole]
    A device that uses exaggeration for emphasis
\end{definition}

\begin{example}
    Some examples of \textbf{Hyperbole}
    
    Her smile was a mile wide. \\
    She broke his heart into a million pieces. \\
    
    Explanation: \\
    
    From this hyperbole, we know that this woman has a very wide grin and is very happy. \\
    The exaggeration \textbf{"a million pieces"} helps us to know just how heartbroken the man is.
\end{example}

\begin{definition}[Understatement]
    A figure of speech in which the words written are less than what a writer or speaker means; the opposite of hyperbole or exaggeration.
\end{definition}

\begin{example}
    Some examples of \textbf{Understatement}
    
    On graduation day, Albert wore five glittering medals around his neck and carried 10 awards under his arm;
    one might say that he had done a decent job during his high school years. \\

    Explanation: \\
    
    With all of these awards and honors, it is clear that Albert did much better than \textbf{"decent."} \\
    This understatement emphasizes just how great he did.
\end{example}

\begin{definition}[Pun]
    A play on words based on the similarity of sound between two words with different meanings.
\end{definition}

\begin{example}
    Some examples of \textbf{Pun}
    
    When the beginning art student accidentally sat down on his classmate's fresh painting, he was sent out of the classroom to change;
    he was not the only one who got a little behind in his work that day. \\

    Explanation: \\
    
    This is a play on the word \textbf{"behind."}
    The student who sat in the wet paint got behind on his work because he had to leave class;
    the student whose painting was sat upon also got a behind in his work—literally.
\end{example}

\begin{definition}[Oxymoron]
    A contradictory combination of words
\end{definition}

\begin{example}
    Some examples of a \textbf{Oxymoron}
    
    \begin{itemize}
        \item A fine mess
        \item Awfully pretty
        \item Serious fun
        \item Working vacation
        \item Sincere lie \\
    \end{itemize}

    Explanation: \\

    \begin{itemize}
        \item A \textbf{mess} is rarely, if ever, fine.
        \item \textbf{"Awful"} usually denotes a negative quality that contradicts \textbf{"pretty."}
        \item \textbf{"Serious"} and \textbf{"fun"} tend to be opposites.
        \item The very idea of vacation is that we leave our work behind.
        \item Being sincere means being honest; a lie is not, by definition, honest.
    \end{itemize}
\end{example}

\begin{definition}[Malapropism]
    Humorous misuse of words—usually by confusing similar-sounding words.
\end{definition}

\begin{example}
    Some examples of \textbf{Malapropism}
    
    \begin{itemize}
        \item The water damage from the storm was so bad that we had to evaporate our home.
        \item Density has brought me to you. \\
    \end{itemize}
    
    Explanation: \\
    
    \begin{itemize}
        \item \textbf{"Evaporate"} and \textbf{"evacuate"} have been confused here.
        \item \textbf{"Density"} is used when \textbf{"destiny"} is what was meant.
    \end{itemize}
\end{example}

\begin{definition}[Verbal Irony]
    A device in which what one means does not match what one says.
\end{definition}

\begin{example}
    Some examples of \textbf{Verbal Irony}
    
    Diana's homecoming date arrived late and left early—without her.
    When her friends took her home, they hugged her and said, \textbf{"What a great night and date!"} \\
    
    Explanation: \\
    
    Diana's friends are trying to jokingly make the point that they know it was a bad night and a terrible date.
    They are using \textbf{verbal irony} because they know it was not a \textbf{"great night"} for their friend.
\end{example}

\begin{definition}[Situational Irony]
    A device in which what one expects to happen does not actually occur.
\end{definition}

\begin{definition}
    \begin{itemize}
        \item A fire station burns to the ground.
        \item I organized my room and cannot find anything now. \\
    \end{itemize}
    
    Explanation: \\
    
    \begin{itemize}
        \item We would not expect a fire station to catch on fire—let alone burn down entirely.
        \item You would expect that getting organized would make it easier to find things, but instead it is more difficult.
    \end{itemize}
\end{definition}

\begin{definition}[Dramatic Irony]
    A device in drama in which the audience knows something that the characters do not.
\end{definition}

\begin{example}
    In a movie or play, a man boards a plane to fly to see his love on the other side of the world; meanwhile,
    the audience knows that his love has already left to fly home and surprise him. \\
    
    Explanation: \\
    
    The audience knows that the man and woman are going to miss each other as they fly to different continents,
    but the characters do not know, making their actions dramatic irony.
\end{example}
    
\newpage